\chapter{End-to-End Algorithm}
\label{ch:end-to-end}

This chapter reassembles every module from Chapters~\ref{ch:sketching}--\ref{ch:scoring} into the
two top-level routines that actually run: \code{map\_reads} (once per program invocation) and
\code{map\_read} (once per read). This is the single most useful reference for a from-scratch
implementation, and is written in its \emph{corrected} form (i.e., incorporating the fixes from
Chapter~\ref{ch:defects} where they affect observable behavior) with explicit call-outs at every
point a decision was made about a defect in the reference source.

\section{\code{map\_reads}: the outer loop}

\begin{algobox}{title={\code{map\_reads}(index, reads\_file, params)}}
\begin{verbatim}
open unmapped_out  := params.params_file + ".unmapped.paf"  if params.params_file is set else None
open paulout       := params.params_file + ".paul.tsv"      if params.params_file is set
                                                                and params.verbose >= 2 else None
buckets := Buckets::new(index)             # allocated once, cleared per read (Section 8.2.1)
run_wide_counters := Counters::new()
run_wide_timers   := Timers::new()

for (query_id, read_seq) in read_fasta(reads_file):
    run_wide_counters.inc("reads")
    per_read_counters, per_read_timers := map_read(index, params, read_seq, query_id,
                                                     buckets, unmapped_out, paulout)
    run_wide_counters += per_read_counters      # accumulate into run totals
    run_wide_timers   += per_read_timers

print_summary_statistics(run_wide_counters, run_wide_timers)
\end{verbatim}
\end{algobox}

\begin{notebox}
Each read is processed with its own fresh, fully-reset per-read counters (see the boxed
correction in \S\ref{sec:map-read-counters} below) that are merged into the run-wide totals
\emph{after} that read finishes --- this merge-after-each-read pattern is what lets a single PAF
line's diagnostic tags (\S\ref{ch:output}) report purely this-read quantities while the
end-of-run summary still reports correct totals across all reads.
\end{notebox}

\section{\code{map\_read}: per-read driver}
\label{sec:map-read}

\subsection{Counter initialization}
\label{sec:map-read-counters}

\begin{bugbox}
The reference C++ does \emph{not} actually reset its per-read counters object between reads: the
call that would clear it (\code{C.clear()}) is commented out, and the subsequent
re-initialization list names only a subset of the counters the function goes on to increment ---
seven counters, including \code{total\_matches} (which feeds directly into the \emph{live}
\code{total\_matches:i:}/\code{match\_inefficiency:f:} PAF tags, not just an end-of-run summary),
are left to accumulate, unreset, across every read a single mapper instance processes, compounding
further every time the (also unreset) running total is merged into the run-wide counters after
each read. The correct, obviously-intended behavior --- and the one specified below --- is to
\textbf{fully reset all per-read counters at the start of every read}, with every counter name the
function goes on to touch pre-registered at zero (so that end-of-read diagnostic reads never fail
on a genuinely-never-incremented counter, e.g.\ \code{mapq60} for a read with a low-confidence
mapping). See Chapter~\ref{ch:defects} for the full analysis and the complete list of affected
counter names.
\end{bugbox}

\begin{algobox}{title={\code{map\_read} --- part 1: sketch, seed, prepare}}
\begin{verbatim}
C := Counters::new()                 # fresh, fully-reset every read (see box above)
T := Timers::new()
buckets.clear()

p := sketcher.sketch(read_seq)                       # Chapter 5
m := len(p)                                            # read's sketch size
C.inc("read_len", len(read_seq))

p_unique := unique_elements_with_info(p)              # Chapter 7; rarest-hit-first

seed_table := {}          # hash -> Seed
diff_hist  := {}          # hash -> remaining occurrence budget
possible_matches := 0
for seed in p_unique:
    seed_table[seed.kmer.h] := seed
    diff_hist[seed.kmer.h]  := seed.occs_in_p
    possible_matches += seed.hits_in_t
C.inc("possible_matches", possible_matches)

lmax   := m
theta  := params.theta
theta2 := theta - params.min_diff
S := one_sweep ? len(p_unique)                        # best-effort -B interpretation, Ch.16
              : floor((1 - theta2) * m) + 1           # Section 7.2

if abs_pos: buckets.set_halflen(len(read_seq))
else:       buckets.set_halflen(m)
if buckets.halflen < MIN_HALFLEN:
    goto unmapped                                      # read too short to bucket meaningfully
\end{verbatim}
\end{algobox}

\begin{algobox}{title={\code{map\_read} --- part 2: match seeds, prune, refine}}
\begin{verbatim}
match_seeds(p_unique, buckets, S)                      # Chapter 8
buckets.propagate_seeds_to_buckets()
sorted_buckets := buckets.get_sorted_buckets()          # most-matches-first
C.inc("seeded_buckets", len(sorted_buckets))

best := match_rest(len(read_seq), m, lmax, p_unique, buckets, sorted_buckets, diff_hist,
                    seed_table, theta, forbidden=None, params.max_overlap, params.metric, params.k)

best2 := None
if best is not None:
    second_best_thr := best.score() * (1 - params.min_diff)
    best2 := match_rest(len(read_seq), m, lmax, p_unique, buckets, sorted_buckets, diff_hist,
                         seed_table, second_best_thr, forbidden=best, params.max_overlap,
                         params.metric, params.k)
\end{verbatim}
\end{algobox}

\subsection{Output: mapped vs.\ unmapped}
\label{sec:map-read-output}

\begin{bugbox}
The reference C++ calls \code{best->set\_global\_stats(...)} and \code{best->print\_paf(...)}
\textbf{unconditionally}, even when \code{best} is empty (every unmapped read) --- dereferencing
an empty \code{std::optional}, undefined behavior in C++ (and simply un-representable against a
proper optional/\code{Option} type in a memory-safe language, since there is no value to
dereference). Per the decision recorded for this document's companion re-implementation, unmapped
reads should instead get a well-defined, minimal record: write \code{query\_id}, read length, and
\code{*}/zero placeholder fields for the alignment-specific columns to the ``unmapped'' output
stream (only opened at all if a params-file/output prefix was configured, matching the
reference's own gating for that stream). See Chapter~\ref{ch:defects} for the full discussion of
why this specific choice (rather than, say, silently skipping unmapped reads entirely) was made.
\end{bugbox}

\begin{algobox}{title={\code{map\_read} --- part 3: emit result}}
\begin{verbatim}
if best is not None:
    if best2 is not None:
        best.set_second_best(best2)              # Chapter 12: feeds MAPQ
    C.inc("matches_in_reported_mappings", best.intersection())
    C.inc("J_best", round(10000 * best.score()))
    C.inc("mappings"); C.inc("mapped_reads")

    best.set_global_stats(theta2, params.min_diff, sketch_size=m, query_id, read_len=len(read_seq),
                           params.k, seeds=S, C.total_matches, C.max_seed_matches, C.seed_matches,
                           C.seeded_buckets, C.final_buckets, fptp=-1.0, map_time=T.secs("query_mapping"))
                            # ^ calc_mapq() runs inside set_global_stats; Chapter 12
    emit_paf_line(stdout, best)                    # Chapter 13
    if best.mapq() == 60: C.inc("mapq60")
    if best.mapq() == 0:  C.inc("mapq0")

    if params.verbose >= 2:
        gt := AnalyseSimulatedReads(query_id, read_seq, diff_hist, m, seed_table, index,
                                     buckets, params.theta)                 # Section 13.x / ground truth
        emit_paf_ground_truth_tags(stdout, gt, best)
        if paulout is not None: gt.print_tsv(paulout)
else:
    emit_minimal_unmapped_record(unmapped_out, query_id, len(read_seq))   # the fix above
    if params.verbose >= 2 and paulout is not None:
        gt := AnalyseSimulatedReads(query_id, read_seq, diff_hist, m, seed_table, index,
                                     buckets, params.theta)
        gt.print_tsv(paulout)

return C, T
\end{verbatim}
\end{algobox}

\section{Full call graph}

\begin{figure}[htbp]
\centering
\begin{tikzpicture}[
    node distance=5mm and 8mm,
    every node/.style={font=\scriptsize},
    n/.style={draw, rounded corners, fill=blue!5, minimum height=7mm, text width=27mm, align=center},
    arr/.style={-{Stealth[length=2mm]}, thin}
  ]
  \node[n] (mr) {map\_read};
  \node[n, below left=of mr] (sketch) {sketch (Ch.5)};
  \node[n, below=of mr] (seed) {unique\_elements\_with\_info (Ch.7)};
  \node[n, below right=of mr] (ms) {match\_seeds (Ch.8)};
  \node[n, below=of ms] (gsb) {get\_sorted\_buckets};
  \node[n, below=of gsb] (mrest) {match\_rest $\times 2$ (best, 2nd-best)};
  \node[n, below left=of mrest] (shp) {seed\_heuristic\_pass (Ch.9)};
  \node[n, below right=of mrest] (fbm) {findBestMapping (Ch.10)};
  \node[n, below=of fbm] (score) {bestFixedLength / lcs};
  \node[n, below=of mrest] (mapq) {calc\_mapq (Ch.12)};
  \node[n, below=of mapq] (out) {emit PAF (Ch.13)};

  \draw[arr] (mr) -- (sketch);
  \draw[arr] (mr) -- (seed);
  \draw[arr] (mr) -- (ms);
  \draw[arr] (ms) -- (gsb);
  \draw[arr] (gsb) -- (mrest);
  \draw[arr] (mrest) -- (shp);
  \draw[arr] (mrest) -- (fbm);
  \draw[arr] (fbm) -- (score);
  \draw[arr] (mrest) -- (mapq);
  \draw[arr] (mapq) -- (out);
\end{tikzpicture}
\caption{Call graph of one \code{map\_read} invocation.}
\end{figure}

\section{Reference C++ implementation (annotated)}

\begin{lstlisting}[style=cppstyle,caption={\code{map\_read}, \texttt{src/shmap.h} --- core structure (elided: verbose-mode/diagnostic branches, identical in spirit to the pseudocode above)}]
inline void map_read(const params_t &params, const FracMinHash &sketcher,
                      ofstream &paulout, ofstream &unmapped_out,
                      const string &query_id, const string &P, BucketsType &B) {
    T.clear();
    B.clear();
    // NOTE: see Chapter 16 -- C.clear() is missing here in the reference source.

    sketch_t p = sketcher.sketch(P);
    qpos_t m = p.size();
    C.inc("read_len", P.size());

    Seeds p_unique = unique_elements_with_info(p);

    h2seed_t p_ht; h2cnt diff_hist; rpos_t possible_matches(0);
    for (const auto &kmer: p_unique) {
        p_ht.insert(make_pair(kmer.kmer.h, kmer));
        diff_hist[kmer.kmer.h] = kmer.occs_in_p;
        possible_matches += kmer.hits_in_T;
    }

    auto lmax = m;
    auto theta = params.theta;
    double theta2 = theta - params.min_diff;
    qpos_t S = qpos_t((1.0 - theta2) * m) + 1;

    if (abs_pos) B.set_halflen(P.size()); else B.set_halflen(m);

    match_seeds(p_unique, B, S);
    B.propagate_seeds_to_buckets();
    auto sorted_buckets = B.get_sorted_buckets();

    std::optional<Mapping> best, best2;
    best = match_rest(P.size(), m, lmax, p_unique, B, sorted_buckets, diff_hist, p_ht,
                       theta, std::nullopt, query_id, &lost_on_pruning, params.max_overlap);
    if (best) {
        double second_best_thr = best->score() * (1.0 - params.min_diff);
        best2 = match_rest(P.size(), m, lmax, p_unique, B, sorted_buckets, diff_hist, p_ht,
                            second_best_thr, best, query_id, &lost_on_pruning, params.max_overlap);
    }

    if (best) {
        if (best2) best->set_second_best(best2.value());
        os = &cout;
    } else {
        os = &unmapped_out;
    }
    // NOTE: see Chapter 16 -- the following two calls are unconditional in the reference
    // source even when `best` is empty; this listing shows the corrected, guarded form.
    best->set_global_stats(theta2, params.min_diff, m, query_id.c_str(), P.size(), params.k,
                            S, C.count("total_matches"), C.count("max_seed_matches"),
                            C.count("seed_matches"), C.count("seeded_buckets"),
                            C.count("final_buckets"), -1.0, T.secs("query_mapping"));
    best->print_paf(*os);
}
\end{lstlisting}
